Un telèfon mòbil conté un motor vibrador que genera un moviment harmònic simple (al llarg de l'eix $y$) al xassís del dispositiu quan es reben trucades o notificacions. Suposem que aquest moviment té una amplitud de 2,00~mm i una freqüència de 150~Hz.

\begin{apartat}{1,25}
Escriviu l'equació de moviment del telèfon si considerem que inicialment es troba a la posició d'equilibri $y = 0\ \mathrm{m}$. Quan es col·loca el telèfon damunt d'una superfície metàl·lica rígida, les vibracions es transmeten al material i generen ones mecàniques transversals que es propaguen a $3\,200\ \mathrm{m\,s^{-1}}$. Calculeu la longitud d'ona de l'ona mecànica que es propaga per la superfície i escriviu l'equació de l'ona corresponent, suposant que la propagació es produeix en el sentit positiu de l'eix $x$, que inicialment es troba en la posició d'equilibri $y = 0\ \mathrm{m}$, i que l'amplitud de l'ona és de 2,00~mm.
\end{apartat}

\begin{apartat}{1,25}
Suposant que el telèfon emet una potència acústica d'1,00~mW durant la vibració i que el so es propaga per l'aire de manera isotròpica (uniformement en totes direccions), calculeu el nivell d'intensitat sonora que percebrà una persona situada a 3,00~m del telèfon.
\end{apartat}

\dades{%
  $I_0 = 10^{-12}\ \mathrm{W\,m^{-2}}$.}
\nota{Assumiu que el front d'ona és esfèric.}
